\section{Empirical Results}
\label{sec:results}

\subsection{Diagnostic Detection Efficacy and Precision (RQ1)}
\label{sec:results:detection}

Table~\ref{tab:detection_results} presents detection validation at the application layer across all eight scenarios.

\begin{table}[htbp]
\centering
\footnotesize
\setlength{\tabcolsep}{3pt}
\caption{Diagnostic Detection Validation (App Layer) against Simulation Ground Truth $I_{\text{comp}}(v)$.}
\label{tab:detection_results}
\begin{tabular}{lrrrrrrrr}
\toprule
\textbf{Scenario} & $n$ & $\rho(Q, I)$ & $F_1$ & \textbf{Top-5} & \textbf{Catalog P} & \textbf{Catalog R} & \textbf{Implicated \%} & \textbf{Gate (s)} \\
\midrule
S01 Autonomous Vehicle & 80 & 0.457 & 0.600 & 0.40 & 0.244 & 0.950 & 97.5\% & 0.59 \\
S02 IoT Smart City & 200 & 0.718 & 0.780 & 0.40 & 0.310 & 0.880 & 71.0\% & 0.56 \\
S03 Financial Trading & 60 & 0.365 & 0.667 & 1.00 & 0.237 & 0.933 & 98.3\% & 0.40 \\
S04 Healthcare & 50 & 0.451 & 0.538 & 0.60 & 0.213 & 0.769 & 94.0\% & 0.16 \\
S05 Hub-and-Spoke & 70 & 0.466 & 0.611 & 0.80 & 0.257 & 1.000 & 100.0\% & 0.92 \\
S06 Microservices Mesh & 90 & 0.405 & 0.565 & 0.40 & 0.261 & 1.000 & 97.8\% & 0.35 \\
S07 Hyper-Scale Enterprise & 300 & 0.480 & 0.627 & 0.20 & 0.251 & 1.000 & 99.7\% & 20.98 \\
S08 Tiny Regression & 12 & 0.538 & 0.333 & 0.60 & 0.250 & 1.000 & 100.0\% & 0.01 \\
\midrule
\textbf{Mean} & --- & \textbf{0.485} & \textbf{0.590} & \textbf{0.550} & \textbf{0.253} & \textbf{0.942} & \textbf{94.8\%} & \textbf{3.00} \\
\bottomrule
\end{tabular}
\end{table}

\paragraph{Key Diagnostic Findings.}
\begin{enumerate}
    \item \textbf{Ranking Efficacy:} Mean Spearman correlation is $\rho = 0.485$ (range $0.365$--$0.718$). Comparing against single-metric baselines, betweenness centrality achieves $\rho = 0.435$, while \textbf{degree centrality alone achieves $\rho = 0.519$}. The value of the RM Diagnostic Pathway lies not in superior scalar ranking, but in \textbf{explainable quality attribution}---routing structural flaws to specific stakeholder roles (Table~\ref{tab:rm_stakeholders}), which a scalar degree count cannot achieve.
    \item \textbf{Catalog Screening Behavior:} The anti-pattern catalog functions as a sensitive screen: mean recall is $0.942$, while precision is $0.253$. The catalog implicates a mean of \textbf{$94.8\%$ of components} at \texttt{CRITICAL}/\texttt{HIGH} severity. This high base-rate confirms that uncalibrated global box-plot thresholds over-flag in dense topologies, underscoring the urgent necessity of delta-aware CI/CD gating (Section~\ref{sec:cicd_gating}).
    \item \textbf{Simpson's Paradox in Node Pooling:} At the system layer (pooling all five node types), the pooled correlation drops to $\rho = 0.028$, while per-type strata show robust positive correlations (Application $0.503$, Broker $0.395$, Node $0.142$). This documents a severe Simpson's paradox resulting from between-type score offsets, proving that stratified evaluation is methodologically essential for heterogeneous software graphs.
\end{enumerate}

\subsection{Prescriptive Refactoring Efficacy (RQ2)}
\label{sec:results:prescriptive}

Table~\ref{tab:prescriptive_results} presents prescriptive refactoring results under per-edit counterfactual verification across all seven benchmark scenarios.

\begin{table}[htbp]
\centering
\small
\setlength{\tabcolsep}{4pt}
\caption{Prescriptive Refactoring Performance under Per-Edit Counterfactual Verification.}
\label{tab:prescriptive_results}
\begin{tabular}{lrrrrrrr}
\toprule
\textbf{Scenario} & \textbf{Base SRI} & \textbf{Mut SRI} & $\Delta\mathrm{SRI}$ & \textbf{Candidates} & \textbf{Admitted} & \textbf{Rejected} & \textbf{Admit \%} \\
\midrule
S01 Autonomous Vehicle & 0.3246 & 0.3160 & +0.0086 & 156 & 45 & 111 & 28.8\% \\
S02 IoT Smart City & 0.3749 & 0.3288 & \textbf{+0.0461} & 391 & 297 & 94 & 76.0\% \\
S03 Financial Trading & 0.3409 & 0.3301 & +0.0108 & 125 & 22 & 103 & 17.6\% \\
S04 Healthcare & 0.3394 & 0.3285 & +0.0109 & 99 & 36 & 63 & 36.4\% \\
S05 Hub-and-Spoke & 0.3223 & 0.3108 & +0.0115 & 127 & 107 & 20 & 84.3\% \\
S06 Microservices Mesh & 0.3293 & 0.3196 & +0.0097 & 163 & 142 & 21 & 87.1\% \\
S07 Hyper-Scale Enterprise & 0.3283 & 0.3112 & \textbf{+0.0171} & 528 & 479 & 49 & 90.7\% \\
\midrule
\textbf{Total / Overall} & --- & --- & --- & \textbf{1589} & \textbf{1128} & \textbf{461} & \textbf{71.0\%} \\
\bottomrule
\end{tabular}
\end{table}

\paragraph{Key Prescriptive Findings.}
\begin{enumerate}
    \item \textbf{Statistically Significant Improvement:} Every scenario achieves a net positive $\Delta\mathrm{SRI}$. A Wilcoxon signed-rank test against zero yields $W=0, p=0.0156 (n=7)$, confirming statistically significant risk reduction across all evaluated domains.
    \item \textbf{Substantial Yield:} Across 1589 generated candidates, \textbf{1128 edits ($71.0\%$)} survive counterfactual verification.
\end{enumerate}

\subsection{Operator Survival and Contributions (RQ3)}
\label{sec:results:operators}

Table~\ref{tab:operator_results} breaks down candidate generation and admission rates by refactoring operator.

\begin{table}[htbp]
\centering
\small
\setlength{\tabcolsep}{5pt}
\caption{Candidate Generation vs. Verified Admission by Refactoring Operator.}
\label{tab:operator_results}
\begin{tabular}{lrrrl}
\toprule
\textbf{Refactoring Operator} & \textbf{Candidates} & \textbf{Admitted} & \textbf{Survival Rate} & \textbf{Primary Structural Target} \\
\midrule
\textbf{Physical Anti-Affinity Reallocation} & 1188 & \textbf{923} & \textbf{77.7\%} & Single points of failure, co-location \\
\textbf{Transport QoS Contract Hardening} & 69 & \textbf{40} & \textbf{58.0\%} & Volatile transport contracts on topics \\
\textbf{Logical Topic Splitting} & 332 & \textbf{165} & \textbf{49.7\%} & High-fan-out message channels, hubs \\
\midrule
\textbf{Total} & \textbf{1589} & \textbf{1128} & \textbf{71.0\%} & --- \\
\bottomrule
\end{tabular}
\end{table}

\textbf{Anti-affinity reallocation demonstrates the highest survival rate ($77.7\%$)}, achieving near-$100\%$ survival in dense topologies (408/409 in Enterprise, 123/123 in Microservices). Transport QoS hardening achieves $58.0\%$ survival where triggered, while topic splitting achieves $49.7\%$.

\subsection{Verification Yield and Compositionality (RQ4)}
\label{sec:results:verification_yield}

\begin{enumerate}
    \item \textbf{Informative Rejection:} All 461 rejected edits record the binding propagation threshold $\theta$ where impact reduction failed to clear $\kappa \sigma_{\text{seed}}$, providing transparent review feedback.
    \item \textbf{Compositional Dynamics:} While every scenario achieves net positive system-level $\Delta\mathrm{SRI}$, two scenarios (S03 Financial Trading and S06 Microservices Mesh) exhibit slight negative mean per-component changes among accepted edits. This demonstrates that individually verified edits can interact sub-additively, validating the necessity of the second-level whole-policy gate.
\end{enumerate}

\subsection{Computational Overhead and CI/CD Feasibility (RQ5)}
\label{sec:results:performance}

\begin{enumerate}
    \item \textbf{Diagnostic Review Gate Latency:} The complete diagnostic pipeline (graph construction, CQP computation, RM attribution, and 18 active anti-pattern detectors) executes in \textbf{0.01~s to 20.98~s} (Table~\ref{tab:detection_results}). The eighteen detectors run in under \textbf{0.2~s total}; computational overhead is dominated by base metric computation.
    \item \textbf{\texttt{DEEP\_PIPELINE} Bounding:} The unconstrained \texttt{DEEP\_PIPELINE} detector attempts exhaustive path enumeration, generating 247,761 findings on tiny fixtures and hanging on larger graphs. Excluding this detector restores sub-second gate execution.
    \item \textbf{Prescriptive Verification Latency:} The full 7-scenario prescriptive sweep (14,301 simulation sweeps) completed in \textbf{2~h 47~min} using multi-core process parallelism (20 workers), establishing its viability for nightly or on-demand CI/CD refactoring runs.
\end{enumerate}
